\chapter{Rational Self-Energy and Schur Nonlinear Eigenproblems}

The third level of the hierarchy begins when converter and auxiliary controller
states are condensed. The reduced problem is no longer a fixed-coefficient QEP.
It is a nonlinear eigenvalue problem whose retained correction is rational in
the spectral parameter. This chapter derives that object and explains why
\(\Pi(s)\), not a static damping matrix, is the general theoretical replacement
after controller or algebraic states are eliminated.

\section{Full and Retained State Partitions}

Let the linearized state matrix be partitioned as
\begin{equation}
  A=
  \begin{bmatrix}
    A_{ee} & A_{ec}\\
    A_{ce} & A_{cc}
  \end{bmatrix}.
  \label{eq:self_full_partition}
\end{equation}
The retained states \(x_e\) may include angles, speeds, or selected interface
variables. The condensed states \(x_c\) may include exciters, governors,
stabilizers, PLLs, voltage and current controllers, plant controllers, filters,
and other auxiliary dynamics. The eigenvalue equation is
\begin{align}
  (sI-A_{ee})x_e-A_{ec}x_c &=0,\\
  -A_{ce}x_e+(sI-A_{cc})x_c &=0.
\end{align}
When \(sI-A_{cc}\) is invertible,
\begin{equation}
  x_c=(sI-A_{cc})^{-1}A_{ce}x_e.
\end{equation}
Substitution gives
\begin{equation}
  S(s)x_e=0,\qquad
  S(s)=sI-A_{ee}-A_{ec}(sI-A_{cc})^{-1}A_{ce}.
  \label{eq:self_schur_operator}
\end{equation}

\begin{proposition}[Exact Schur equivalence]
If \(sI-A_{cc}\) is invertible, then \(s\) is an eigenvalue of the full matrix
\(A\) with retained component \(x_e\ne0\) if and only if \(S(s)x_e=0\).
\end{proposition}

\begin{proof}
The proof follows by the two substitutions above. Starting from a full
eigenpair gives \(S(s)x_e=0\). Conversely, if \(S(s)x_e=0\), define
\(x_c=(sI-A_{cc})^{-1}A_{ce}x_e\). Then the block eigenvalue equations are
satisfied. \qed
\end{proof}

The result is exact, but the retained operator is nonlinear in \(s\). Any
constant reduced matrix that tries to replace \(A_{ec}(sI-A_{cc})^{-1}A_{ce}\)
must be an approximation.

\section{Descriptor Schur Complement}

Power-system linearizations are often descriptor systems rather than explicit
ODEs. The more general eigenvalue equation is
\begin{equation}
  (sE-A)x=0.
  \label{eq:self_descriptor_pencil}
\end{equation}
Partition retained and condensed variables as
\begin{equation}
  sE-A
  =
  \begin{bmatrix}
    sE_{rr}-A_{rr} & sE_{rc}-A_{rc}\\
    sE_{cr}-A_{cr} & sE_{cc}-A_{cc}
  \end{bmatrix}.
  \label{eq:self_descriptor_partition}
\end{equation}
When \(sE_{cc}-A_{cc}\) is invertible, the exact retained pencil is
\begin{equation}
  S(s)
  =
  sE_{rr}-A_{rr}
  -
  (sE_{rc}-A_{rc})
  (sE_{cc}-A_{cc})^{-1}
  (sE_{cr}-A_{cr}).
  \label{eq:self_descriptor_schur}
\end{equation}
This form is the correct starting point for DAEs, semi-explicit models,
algebraic variables that are not eliminated first, and identified port models.
The explicit-state formula \eqref{eq:self_schur_operator} is the special case
\(E=I\) with compatible zero off-diagonal descriptor blocks. The descriptor
form also makes visible which eliminated poles are genuine dynamic poles and
which are algebraic or infinite-eigenvalue artifacts of the chosen
representation.

\section{Port Equivalence of Self-Energy}

The self-energy must be treated as a port map, not as a property of one hidden
state realization. Let two condensed subsystems have hidden coordinates
\(z_1\) and \(z_2\), possibly of different dimensions, and suppose they produce
the same retained correction on a spectral domain \(\mathcal D\):
\begin{equation}
  \Pi_1(s)=\Pi_2(s),\qquad s\in\mathcal D .
  \label{eq:self_port_equiv_pi}
\end{equation}
Then every retained-port statement that uses only
\begin{equation}
  T(s)=s^2M+sD_0+L+\Pi(s)
  \label{eq:self_port_equiv_T}
\end{equation}
is invariant under the change of hidden realization.

\begin{theorem}[Spectral sufficiency of the port self-energy]
\label{thm:self_port_sufficiency}
Assume that the retained matrices \(M,D_0,L\) are the same and that two hidden
subsystems satisfy \eqref{eq:self_port_equiv_pi} on a domain \(\mathcal D\)
that contains the protected roots and the contours used for certification.
Then the two reduced operators have the same retained roots in
\(\mathcal D\), the same retained Green function \(T(s)^{-1}\) wherever the
inverse exists, and the same low-rank action determinants for any action that
is expressed only through the retained ports.
\end{theorem}

\begin{proof}
Under \eqref{eq:self_port_equiv_pi}, the two retained operators
\eqref{eq:self_port_equiv_T} are identical on \(\mathcal D\). Equality of
roots, inverses on the resolvent set, and any determinant formula built from
that inverse follows directly. The full hidden-state eigenvectors may differ,
but those coordinates are not observables of a retained-port decision.
\end{proof}

This theorem is deliberately modest. It does not say that hidden
participation factors, Euclidean controller-state norms, or internal energy
coordinates are invariant. It says that a retained-port planning claim should
depend on the map \(\Pi(s)\), the port normalization, and the protected
spectral region. Similarity transformations \(z=S\widehat z\), changes of
hidden units, or a different order of Schur elimination are harmless only to
the extent that they preserve the same port self-energy on the domain being
used.

\section{General Second-Order Self-Energy Form}

The exact Schur operator \eqref{eq:self_schur_operator} does not by itself imply
that the reduced second-order model contains only a dynamic damping term. To see
the required structure, write a retained mechanical model with eliminated
controller states \(z\):
\begin{align}
  M\ddot\theta+D_0\dot\theta+\Lt\theta+Gz &=0,\\
  \dot z &= A_{cc}z+B_\theta\theta+B_\omega\dot\theta .
  \label{eq:self_structured_controller}
\end{align}
After Laplace transformation and elimination of \(z\),
\begin{equation}
  T(s)
  =
  s^2M+sD_0+\Lt+\Pi(s),
  \qquad
  \Pi(s)=G(sI-A_{cc})^{-1}(B_\theta+sB_\omega).
  \label{eq:self_general_pi}
\end{equation}
The sign of \(G,B_\theta,B_\omega\) depends on the port convention; the
important point is the channel separation. In general, \(\Pi(s)\) contains a
rational stiffness contribution and a rational damping contribution:
\begin{equation}
  \Pi(s)=\Pi_K(s)+s\Pi_D(s).
  \label{eq:self_pi_decomposition}
\end{equation}
Only under the additional assumption \(\Pi_K(s)=0\), or after a defensible
absorption of \(\Pi_K\) into a stated stiffness correction, does the operator
reduce to the dynamic-damping form
\begin{equation}
  T(s)=s^2M+s\Sigma(s)+\Lt,
  \label{eq:self_second_order}
\end{equation}
where
\begin{equation}
  \Sigma(s)=D_0+\Pi_D(s)
  =
  D_0+C(sI-A_{cc})^{-1}B
  \quad\hbox{in the damping-channel special case.}
  \label{eq:self_sigma_def}
\end{equation}
This is the precise transition from a fixed \(D\) to a rational
self-energy \(\Pi(s)\). The broader self-energy is \(\Pi(s)\);
\(\Sigma(s)\) names the dynamic damping component.

The notation \(\Sigma(s)\) is therefore deliberately restricted. It is not
simply damping in the classical viscous sense. It can include dissipative,
anti-dissipative, phase-lag, resonant, and cross-channel effects, depending on
the poles and residues of the condensed block. Its Hermitian part on
\(s=j\omega\) may suggest dissipative behavior over a frequency band, but the
full eigenvalue effect also depends on eigenvectors and on \(T_s(s)\).

\section{Dynamic Green Function and Effective Resistance}

The inverse of the retained operator is the dynamic Green function of the
condensed network:
\begin{equation}
  G_T(s)=T(s)^{-1},
  \qquad s\in\rho(T),
  \label{eq:self_green_function}
\end{equation}
where \(\rho(T)\) denotes the resolvent set. This object is the map from a
retained port injection to the retained response at spectral parameter \(s\).
It is therefore the natural carrier of dynamic grid strength: it contains the
graph stiffness, inertia, fixed damping, controller self-energy, and
non-normal amplification in one operator.

\begin{definition}[Dynamic effective resistance]
For a line or candidate line \(e=(i,j)\), let \(a_e=e_i-e_j\). The dynamic
effective resistance of the line is
\begin{equation}
  \mathcal R_e(s)=a_e^TG_T(s)a_e
  =
  a_e^T
  \left[
    s^2M+sD_0+L+\Pi(s)
  \right]^{-1}
  a_e .
  \label{eq:self_dynamic_resistance}
\end{equation}
\end{definition}

When \(s=0\), \(\Pi(0)=0\), and the uniform angle reference is removed, this
quantity reduces to the classical graph effective resistance
\begin{equation}
  \mathcal R_e(0)=a_e^TL^\dagger a_e .
  \label{eq:self_static_resistance_limit}
\end{equation}
For \(s\ne0\), \(\mathcal R_e(s)\) is not a physical resistor. It is a
meromorphic port compliance. Its poles are the retained roots of \(T(s)\), and
its residues measure how strongly each mode is seen across the endpoint
difference of line \(e\).

\begin{theorem}[Exact finite line-action law]
\label{thm:self_finite_line_action}
Let a finite line-weight change enter the retained operator as
\begin{equation}
  T_{\rm new}(s)=T(s)+\Delta b_ea_ea_e^T .
  \label{eq:self_line_action_update}
\end{equation}
For all \(s\) where \(T(s)\) is invertible,
\begin{equation}
  \det T_{\rm new}(s)
  =
  \det T(s)\left[1+\Delta b_e\mathcal R_e(s)\right].
  \label{eq:self_line_action_determinant}
\end{equation}
Thus any new or shifted root created by this rank-one line action and not
already a root of \(T\) satisfies
\begin{equation}
  1+\Delta b_e\mathcal R_e(s)=0 .
  \label{eq:self_line_action_scalar}
\end{equation}
\end{theorem}

\begin{proof}
Apply the matrix determinant lemma to
\(T(s)+\Delta b_ea_ea_e^T\). This gives
\[
  \det(T+\Delta b_ea_ea_e^T)
  =
  \det T\left(1+\Delta b_ea_e^TT^{-1}a_e\right),
\]
which is \eqref{eq:self_line_action_determinant}.
\end{proof}

This result is stronger than a local derivative. It gives an exact scalar
equation for a finite rank-one line modification even when \(T(s)\) is a large
rational NEP. The same equation also explains why classical effective
resistance is a static shadow of a richer dynamic object: the graph quantity
is the zero-frequency limit of the port Green function.

\section{Low-Rank Dynamic Action Calculus}

Line reinforcement is only one member of a broader action family. Suppose a
planning or retuning action enters through a low-rank retained-port update
\begin{equation}
  T_u(s)=T(s)+U(s)K_u(s)V(s)^* .
  \label{eq:self_low_rank_action}
\end{equation}
Define the action interaction matrix
\begin{equation}
  \mathcal G_u(s)=V(s)^*T(s)^{-1}U(s).
  \label{eq:self_action_green_matrix}
\end{equation}
Then the determinant lemma gives
\begin{equation}
  \det T_u(s)
  =
  \det T(s)\det\!\left[I+K_u(s)\mathcal G_u(s)\right].
  \label{eq:self_low_rank_determinant}
\end{equation}
The large NEP has been reduced, for the purpose of this action, to a
determinant whose dimension is the number of action ports.

Several planning levers become special cases of
\eqref{eq:self_low_rank_determinant}:
\begin{align}
  \hbox{line } e:\quad
  &U=V=a_e,\qquad K_u(s)=\Delta b_e, \label{eq:self_action_line}\\
  \hbox{local damping at bus } i:\quad
  &U=V=e_i,\qquad K_u(s)=s\Delta d_i, \label{eq:self_action_damping}\\
  \hbox{local inertia at bus } i:\quad
  &U=V=e_i,\qquad K_u(s)=s^2\Delta m_i, \label{eq:self_action_inertia}\\
  \hbox{controller retuning through ports}:\quad
  &U=U_c,\quad V=V_c,\quad K_u(s)=\Delta K_c(s).
  \label{eq:self_action_controller}
\end{align}
Thus line, damping, inertia, and controller actions are not four unrelated
sensitivity formulas. They are four kernels projected through the same Green
function \(T(s)^{-1}\).

\section{Modal Sensitivity as a Local Limit}

The usual NEP sensitivity is recovered by expanding the Green function near a
simple root. Let \(T(s_k)x_k=0\), \(y_k^*T(s_k)=0\), and
\(y_k^*T_s(s_k)x_k\ne0\). Then
\begin{equation}
  T(s)^{-1}
  \sim
  \frac{x_ky_k^*}
  {(s-s_k)y_k^*T_s(s_k)x_k}.
  \label{eq:self_green_pole_expansion}
\end{equation}
For a line \(e\),
\begin{equation}
  \mathcal R_e(s)
  \sim
  \frac{(a_e^Tx_k)(y_k^*a_e)}
  {(s-s_k)y_k^*T_s(s_k)x_k}.
  \label{eq:self_resistance_residue}
\end{equation}
Substituting this expansion into the finite-action equation
\eqref{eq:self_line_action_scalar} gives
\begin{equation}
  s-s_k
  \approx
  -\Delta b_e
  \frac{y_k^*a_ea_e^Tx_k}
  {y_k^*T_s(s_k)x_k},
  \label{eq:self_line_local_shift}
\end{equation}
which is exactly the first-order line sensitivity
\begin{equation}
  \frac{ds_k}{db_e}
  =
  -
  \frac{y_k^*a_ea_e^Tx_k}
  {y_k^*T_s(s_k)x_k}.
  \label{eq:self_line_sensitivity_limit}
\end{equation}
The sensitivity is therefore not an isolated formula. It is the local residue
limit of an exact finite-action law.

\section{Pole-Residue Structure}

Assume for exposition that \(A_{cc}\) is diagonalizable:
\begin{equation}
  A_{cc}=R\Lambda_c R^{-1},\qquad
  \Lambda_c=\diag(\alpha_1,\ldots,\alpha_m).
\end{equation}
Then
\begin{equation}
  C(sI-A_{cc})^{-1}B
  =
  \sum_{\ell=1}^m
  \frac{C r_\ell \ell_\ell^T B}{s-\alpha_\ell},
  \label{eq:self_partial_fraction}
\end{equation}
where \(r_\ell\) and \(\ell_\ell\) are right and left eigenvectors of
\(A_{cc}\), normalized so that \(\ell_\ell^Tr_\ell=1\). The residue
\(C r_\ell \ell_\ell^T B\) measures how strongly the eliminated controller pole
\(\alpha_\ell\) is seen by retained coordinates.

This expression explains why a nearby controller pole can dominate a damping
margin. If \(s\) is close to \(\alpha_\ell\), the denominator is small. If the
residue projects strongly onto the retained eigenvectors, the self-energy term
can dominate the eigenvalue sensitivity.

\section{NEP Sensitivity with Self-Energy}

For \(T(s,\rho)v=0\), the simple-root sensitivity is
\begin{equation}
  \frac{ds_\star}{d\rho}
  =
  -\frac{u^H T_\rho(s_\star,\rho)v}
  {u^H T_s(s_\star,\rho)v}.
  \label{eq:self_nep_sens}
\end{equation}
For the general operator \eqref{eq:self_general_pi},
\begin{equation}
  T_s(s)=2sM+D_0+\Pi_s(s).
  \label{eq:self_Ts_general}
\end{equation}
In the damping-channel specialization \eqref{eq:self_second_order},
\begin{equation}
  T_s(s)=2sM+\Sigma(s)+s\Sigma'(s),
  \label{eq:self_Ts}
\end{equation}
with
\begin{equation}
  \Sigma'(s)=-C(sI-A_{cc})^{-2}B.
  \label{eq:self_sigma_prime}
\end{equation}
Thus even the denominator of the sensitivity contains controller-resolvent
terms. A fixed scalar damping bridge cannot reproduce this dependence unless
the rational term is negligible over the critical spectral region.

The additional denominator term \(s\Sigma'(s)\) is the precise mathematical
place where the condensed controller block retains spectral memory. It is not
another static damping coefficient; it is a slope of the self-energy with
respect to the eigenvalue being solved for. Using
\eqref{eq:self_sigma_prime}, this slope contains the squared resolvent
\((sI-A_{cc})^{-2}\). Hence it becomes large when the retained pole is close to
a visible condensed-controller pole. In that regime, replacing
\(\Sigma(s)\) by a constant matrix may preserve one fitted value of the
self-energy while losing the local pole-motion geometry.

\section{Why Static Damping Bridges Fail}

A static bridge assumes that the reduced damping mechanism can be represented
by a constant matrix \(D_{\rm red}\). In the self-energy theory, that assumption
requires
\begin{equation}
  C(sI-A_{cc})^{-1}B\approx D_{\rm red}-D_0
  \label{eq:self_static_approx}
\end{equation}
over all poles and frequency bands relevant to the claim. This is a strong
condition. It can hold locally if the controller poles are far away and the
resolvent varies slowly. It fails when the critical pole is near a controller
pole, when residues are large, or when the relevant mode is dominated by
condensed states.

The failure is structural rather than cosmetic. A constant matrix has no poles.
It cannot represent the denominator \(s-\alpha_\ell\). It can approximate the
rational term at a selected point, but it cannot preserve the full spectral
mechanism over a region unless additional assumptions are imposed.

\begin{remark}[Static models are local models]
The conclusion is not that every static reduced model is useless. The
conclusion is that a static bridge has a validity region. It is defensible
when controller poles are far from the protected region, when residues are
small in the retained eigenvector directions, and when \(\Sigma'(s)\) is small
enough that the denominator of \eqref{eq:self_nep_sens} is not materially
changed. It is not defensible as a universal damping-margin surrogate near a
controller resonance.
\end{remark}

\section{Certified Adequacy of Reduced Operators}

The hierarchy is useful only if it contains an escalation rule. Let
\(\widehat T(s)\) be a reduced or approximate operator, such as a static bridge,
a fitted self-energy, or a lower-order rational approximation of \(\Pi(s)\).
On a closed contour \(\Gamma\) that bounds a protected region, define
\begin{equation}
  \eta_\Gamma
  =
  \sup_{s\in\Gamma}
  \left\|
  \widehat T(s)^{-1}\left[T(s)-\widehat T(s)\right]
  \right\|_2 .
  \label{eq:self_eta_gamma}
\end{equation}

\begin{theorem}[Contour adequacy certificate]
\label{thm:self_contour_adequacy}
Assume that \(T\) and \(\widehat T\) are regular analytic matrix functions on
and inside \(\Gamma\), that \(\widehat T(s)\) is nonsingular on \(\Gamma\), and
that
\begin{equation}
  \eta_\Gamma<1 .
  \label{eq:self_eta_condition}
\end{equation}
Then \(T\) and \(\widehat T\) have the same number of eigenvalues inside
\(\Gamma\), counted with algebraic multiplicity.
\end{theorem}

\begin{proof}
The statement is the matrix-function analogue of Rouche's theorem. Since
\(\|\widehat T^{-1}(T-\widehat T)\|_2<1\) on \(\Gamma\), the homotopy
\(\widehat T+\tau(T-\widehat T)\) remains nonsingular on \(\Gamma\) for
\(0\le\tau\le1\). The argument principle for regular analytic matrix
functions then preserves the eigenvalue count inside the contour
\cite{kato1995perturbation,guettel2017nep}.
\end{proof}

This certificate turns the operator hierarchy into a computational policy:
\begin{equation}
  \eta_\Gamma<1
  \quad\Longrightarrow\quad
  \hbox{the reduced operator preserves the protected eigenvalue count,}
  \label{eq:self_eta_policy_accept}
\end{equation}
whereas
\begin{equation}
  \eta_\Gamma\ge1
  \quad\Longrightarrow\quad
  \hbox{escalate to a richer NEP, identified self-energy, or full DAE.}
  \label{eq:self_eta_policy_escalate}
\end{equation}

A local decision also needs a sign certificate. If an approximate model
predicts a margin derivative \(\widehat g\), and an error analysis gives
\begin{equation}
  |g-\widehat g|\le\varepsilon_g ,
  \label{eq:self_sign_error_bound}
\end{equation}
then
\begin{equation}
  |\widehat g|>\varepsilon_g
  \quad\Longrightarrow\quad
  \operatorname{sign}(g)=\operatorname{sign}(\widehat g).
  \label{eq:self_sign_certificate}
\end{equation}
Without a bound of this type, a ranked action is a screen rather than a
certified recommendation.

\section{Data-Identified Self-Energy}

The self-energy theory should not require access to proprietary hidden
matrices. If a manufacturer, EMT model, or field experiment provides only
terminal data, the retained-port object can be estimated as
\begin{equation}
  \widehat\Pi(s)
  \approx
  \Pi(s)
  \label{eq:self_identified_pi}
\end{equation}
from frequency scans, injected perturbations, PMU records, or black-box
simulation. The estimation method may be vector fitting, Loewner realization,
gray-box calibration, or a structure-preserving DAE identification method.
Recent converter-identification work is moving in this direction by combining
terminal time-series data with physics-informed converter structure and
frequency-domain validation \cite{darii2026graybox}; recent
port-Hamiltonian DAE identification similarly emphasizes preserving
interconnection and passivity structure under noisy input-output data
\cite{hagelaars2026identification}.

An identified self-energy is useful for planning only if it comes with an
error model. On a protected frequency or contour set \(\Omega\), require
\begin{equation}
  \|\Pi(s)-\widehat\Pi(s)\|_2\le\varepsilon_\Pi(s),
  \qquad s\in\Omega .
  \label{eq:self_pi_identification_error}
\end{equation}
This bound propagates through the same adequacy certificate:
\begin{equation}
  \widehat T(s)=s^2M+sD_0+L+\widehat\Pi(s),
  \qquad
  \eta_\Gamma
  =
  \sup_{s\in\Gamma}
  \left\|
  \widehat T(s)^{-1}\left[\Pi(s)-\widehat\Pi(s)\right]
  \right\|_2 .
  \label{eq:self_identified_eta}
\end{equation}
The practical chain is therefore
\begin{equation}
  \hbox{data}
  \rightarrow
  \widehat\Pi(s)
  \rightarrow
  \varepsilon_\Pi
  \rightarrow
  \eta_\Gamma
  \rightarrow
  \hbox{certified or rejected action}.
  \label{eq:self_data_to_certification}
\end{equation}
The thesis contribution is not merely to identify a converter. It is to decide
how much identification error can be tolerated before a planning recommendation
changes sign or loses its protected eigenvalue count.

\section{Mode Families}

The Schur partition motivates a mode-family diagnostic, but a raw Euclidean
ratio of retained and condensed state norms is not invariant under a change of
units or controller realization. If physically meaningful weights are available,
use the weighted index
\begin{equation}
  \pi_c^{(W)}
  =
  \frac{x_c^H W_c x_c}
  {x_e^H W_e x_e+x_c^H W_c x_c},
  \label{eq:self_control_participation}
\end{equation}
where \(W_e,W_c\succeq0\) represent compatible energy, power, or per-unit
normalization choices. Small \(\pi_c^{(W)}\) indicates a retained-state or
network-family mode. Large \(\pi_c^{(W)}\) indicates a control-family mode.
Intermediate values indicate mixed modes.

When no defensible \(W_c\) exists, the more invariant diagnostic is not internal
state size but port visibility:
\begin{equation}
  \mathcal V_c(s)
  =
  \norm{C(sI-A_{cc})^{-1}B}_2 .
  \label{eq:self_controller_visibility}
\end{equation}
This quantity asks whether the condensed block is visible from retained
coordinates at the frequency being tested. It changes under input-output
scaling, so the ports must still be normalized, but it is not changed by a
similarity transformation of the hidden controller state.

\section{Structured Pseudospectrum and Fragility Radius}

Nominal eigenvalues do not fully describe a non-normal converter-rich model.
The unstructured NEP pseudospectrum is
\begin{equation}
  \Lambda_\varepsilon(T)
  =
  \left\{
  s:\sigma_{\min}(T(s))\le\varepsilon
  \right\}.
  \label{eq:self_unstructured_pseudospectrum}
\end{equation}
This set is useful as a numerical warning, but it permits arbitrary matrix
perturbations. Planning perturbations are not arbitrary: they come from line
uncertainty, damping placement, inertia placement, controller gains, delays,
operating-point errors, or identified-model errors. The relevant object is
therefore the action-structured pseudospectrum
\begin{equation}
  \Lambda_\varepsilon^{\mathcal A}(T)
  =
  \left\{
  s:\exists u\in\mathcal A,\ 
  \|W_uu\|\le\varepsilon,\ 
  \det T(s,u)=0
  \right\},
  \label{eq:self_structured_pseudospectrum}
\end{equation}
where \(\mathcal A\) is the declared physical action or uncertainty family and
\(W_u\) normalizes units and costs. Recent nonlinear spectral algorithms make
the broader computation of spectra and pseudospectra increasingly relevant for
NEPs, but the power-system decision must still restrict perturbations to
physically admissible directions \cite{colbrook2025universal}.

For a protected spectral boundary \(\partial\mathcal S\), define the dynamic
fragility radius
\begin{equation}
  r_{\rm frag}
  =
  \inf
  \left\{
  \|W_uu\|:
  \exists s\in\partial\mathcal S,\ 
  \det T(s,u)=0
  \right\}.
  \label{eq:self_fragility_radius}
\end{equation}
This number answers the planning question directly: what is the smallest
physical perturbation that can violate the protected margin? A damping-ratio
margin, an abscissa margin, a distance-to-resonance margin, a passivity margin,
or a protected-gain margin can each be used to define
\(\partial\mathcal S\). Local pole-motion cones and hidden-margin derivatives
are then first-order approximations to \eqref{eq:self_fragility_radius}, not
substitutes for it.

\section{Passivity and Resolvent Gates}

A stable \(A_{cc}\) block does not guarantee that the condensed self-energy is
benign. A passivity statement requires a declared input-output pairing and a
supply rate in the sense of dissipativity theory \cite{willems1972dissipative}.
For a MIMO retained port, reducing passivity to one scalar index can hide the
direction that causes the deficiency. With input and output normalizers
\(W_u\) and \(W_y\), define the directional passivity matrix
\begin{equation}
  \mathcal P_\Sigma(\omega)
  =
  \frac12
  \left[
  W_y\Sigma(j\omega)W_u^{-1}
  +
  W_u^{-*}\Sigma(j\omega)^*W_y^*
  \right].
  \label{eq:self_directional_passivity}
\end{equation}
This matrix records both intensity and direction of passivity excess or
deficit. Recent matrix-valued passivity work is useful here precisely because
scalar indices can lose MIMO channel structure \cite{ru2026matrix}. The
unweighted Hermitian screen used below is the special case
\begin{equation}
  H_\Sigma(\omega)=
  \frac{\Sigma(j\omega)+\Sigma(j\omega)^H}{2}.
  \label{eq:self_hermitian}
\end{equation}
Negative eigenvalues of \(H_\Sigma(\omega)\), or of the weighted matrix
\(\mathcal P_\Sigma(\omega)\), indicate frequencies where the condensed
contribution acts like negative damping in a retained port direction. Large
singular values of the resolvent identify frequencies where the eliminated
block strongly amplifies retained-state motion. Dynamic passivity multipliers
provide a broader certificate family when strict passivity in the trivial
metric is too conservative \cite{gorbunov2026dynamic}.

Recent passivity-based grid-forming work suggests a useful interpretation of
this gate \cite{he2024passivity,gui2024passivity}. If a controller parameter
\(\kappa\) changes the condensed block, write
\(A_{cc}=A_{cc}(\kappa)\) and define the band-limited self-energy passivity
screen
\begin{equation}
  \mu_\Sigma(\Omega;\kappa)
  =
  \inf_{\omega\in\Omega}
  \lambda_{\min}\!\left(H_\Sigma(\omega;\kappa)\right),
  \label{eq:self_passivity_index}
\end{equation}
over the frequency band \(\Omega\) that contains the protected electromechanical
or controller-interaction modes. If \(\mu_\Sigma(\Omega;\kappa)>0\), the condensed
controller contribution is dissipative in the retained coordinates over that
band. If \(\mu_\Sigma(\Omega;\kappa)<0\), there exists a retained direction and a
frequency where the condensed block supplies negative damping.

\begin{theorem}[Local self-energy passivity-retuning certificate]
\label{thm:self_passivity_retuning}
Let \(\Omega\subset\R\) be compact. Assume \(A_{cc}(\kappa)\), \(B(\kappa)\),
and \(C(\kappa)\) are continuously differentiable on an interval
\(\mathcal K\), and that \(j\omega I-A_{cc}(\kappa)\) is invertible for all
\((\omega,\kappa)\in\Omega\times\mathcal K\). At a fixed \(\kappa\), suppose
the infimum in \eqref{eq:self_passivity_index} is attained at a unique
\(\omega_\star\), and that the smallest eigenvalue of
\(H_\Sigma(\omega_\star;\kappa)\) is simple with unit eigenvector \(q_\star\).
Then \(\mu_\Sigma\) is differentiable at \(\kappa\), and
\begin{equation}
  \frac{d\mu_\Sigma}{d\kappa}
  =
  q_\star^H
  \frac{\partial H_\Sigma(\omega_\star;\kappa)}{\partial\kappa}
  q_\star .
  \label{eq:self_passivity_envelope}
\end{equation}
If this quantity is nonnegative throughout an interval, then
\(\mu_\Sigma\) is nondecreasing on that interval. If it is positive before a
single zero and negative afterward, the band-limited passivity margin is locally
unimodal and the zero gives a candidate retuning optimum.
\end{theorem}

\begin{proof}
For fixed \(\omega\), \(H_\Sigma(\omega;\kappa)\) is Hermitian and differentiable
under the stated invertibility assumptions. The standard simple-eigenvalue
perturbation formula gives
\(\partial_\kappa\lambda_{\min}=q_\star^H(\partial_\kappa H_\Sigma)q_\star\)
at the active frequency. Since the margin is an infimum over the compact
frequency set, Danskin's envelope theorem applies to the active minimizer
\cite{danskin1967maxmin,rockafellar1998variational}. Under uniqueness and
simplicity, the directional formula collapses to
\eqref{eq:self_passivity_envelope}. Monotonicity follows by integrating the
derivative over the interval. The unimodal statement is the same first-order
condition with one sign change.
\end{proof}

The derivative in \eqref{eq:self_passivity_envelope} is computable from the same
resolvent that defines the self-energy. With \(R(s,\kappa)=(sI-A_{cc}(\kappa))^{-1}\),
\begin{equation}
  \frac{\partial\Sigma}{\partial\kappa}
  =
  C_\kappa R B
  +
  C R_\kappa B
  +
  C R B_\kappa,\qquad
  R_\kappa
  =
  R\,\frac{\partial A_{cc}}{\partial\kappa}\,R.
  \label{eq:self_sigma_kappa_derivative}
\end{equation}
If only the controller matrix \(A_{cc}\) depends on \(\kappa\), the first and
third terms vanish. The sign of \(d\mu_\Sigma/d\kappa\) is therefore not a
generic algebraic fact; it is a physical property of the controller family,
operating point, and frequency band.

For a grid-following PLL, this distinction matters. Increasing PLL bandwidth can
improve phase tracking over one range and degrade weak-grid damping over another.
The thesis therefore treats monotonicity as a local certificate, not as a global
law. The actionable design problem is
\begin{equation}
  \kappa_\star\in
  \operatorname*{arg\,max}_{\kappa\in\mathcal K}
  \mu_\Sigma(\Omega;\kappa),
  \qquad
  \text{subject to a protected-pole gate.}
  \label{eq:self_passivity_optimal_kappa}
\end{equation}
This formulation also explains what must be validated after retuning:
\(\mu_\Sigma\) should increase in the relevant band, and the protected pole
should move in the desired direction.

This is not a complete stability theorem for the retained NEP. It is a
certificate candidate. The missing companion step is the pole sensitivity:
\begin{equation}
  \frac{d}{d\rho}\mu_\Sigma(\Omega;\rho),
  \qquad
  \frac{ds_\star}{d\rho}
  =
  -\frac{u^H T_\rho(s_\star,\rho)v}
  {u^H T_s(s_\star,\rho)v}.
  \label{eq:self_passivity_retuning}
\end{equation}
A controller-retuning action is credible only if it increases the passivity
screen in the relevant band and moves the protected pole in the correct
direction. This gives a principled version of ``PLL retuning'': the action is
not merely moving an eigenvalue numerically; it is attempting to restore a
band-limited dissipativity margin of the Schur self-energy.

These diagnostics are gates, not standalone stability theorems. They must be
connected to eigenvalue sensitivities and full-model validation before they can
support planning decisions. Their theoretical value is that they point to the
operator that must be inspected: \(\Sigma(s)\), not a scalar damping constant.

\section{Sign Characteristic and Controller-Resolvent Collisions}

The language of Krein signature is mature for Hamiltonian and self-adjoint
spectral problems, and the matrix-polynomial literature uses the more precise
term \emph{sign characteristic} \cite{gohberg1982matrixpolynomials,
gohberg1983meromorphic,mehrmann2016sign}. Kollár and Miller give a useful
graphical interpretation through Evans--Krein functions for polynomial operator
pencils \cite{kollar2014graphical}. That literature should not be cited as if it
already proved the converter result needed here. The thesis uses it as the
correct mathematical foundation for a narrower question: can a rational,
controller-condensed self-energy produce an operational sign collision that
predicts weak-grid synchronization loss?

For a Hermitian matrix polynomial, the sign characteristic is invariant under
structure-preserving transformations and is tied to the perturbation behavior of
real or imaginary-axis eigenvalues. A converter-condensed operator
\[
  T(s)=s^2M+sD_0+\Lt+\Pi(s)
\]
is usually not Hermitian as a function of \(s\), and its rational term may be
dissipative and non-normal. Therefore the following object is only an
operational signature gate, not a completed Krein theorem:
\begin{equation}
  \chi_k^{\rm op}
  =
  \operatorname{sign}
  \Real\!\left(v_k^H T_s(s_k)v_k\right),
  \qquad
  T(s_k)v_k=0.
  \label{eq:self_operational_signature}
\end{equation}
If left and right eigenvectors differ strongly, the bilinear variant
\(u_k^H T_s(s_k)v_k\) should be tracked together with eigenvalue conditioning.
If \(s_k=j\omega_k\) lies on a symmetry line where a Hermitian reduction is
available, one may instead inspect the inertia of
\((T_s(j\omega_k)+T_s(j\omega_k)^H)/2\) on the active eigenspace.

The reason this is relevant to controller-resonant instability is the derivative
term
\[
  T_s(s)=2sM+D_0+\Pi_s(s)
  \quad
  \hbox{or, in the damping-channel case,}\quad
  T_s(s)=2sM+\Sigma(s)+s\Sigma'(s).
\]
Near a pole of \(A_{cc}\), the resolvent derivative contains squared-resolvent
growth. A mode can then become dominated by the controller self-energy rather
than by the mechanical part of the QEP. The proposed mechanism is a
controller-resolvent sign collision: two tracked modes approach while their
operational signatures differ, and the denominator
\(u^H T_s(s)v\) in the NEP sensitivity becomes small or ill-conditioned.

This bridge is a research claim, not a settled theorem. It is credible only if a
benchmark shows the same event in three views: a sign change or singularity in
\eqref{eq:self_operational_signature}, growth of the absolute self-energy
dominance index, and a full-model eigenvalue or time-domain loss of damping. The
general theory of dissipation-induced instabilities explains why such collisions
are plausible \cite{krechetnikov2007dissipation}; the new power-system claim
would be that the Schur self-energy identifies the controller channel that
causes them.
